% !TEX root = main.tex
\section{System and Threat Model}
\label{sec:system_threat}

We consider embedded systems based on FPGA SoC platforms, ARM Cortex-A processors with Armv9.2-A RME, and ARM Cortex-M microcontroller units (MCUs).
FPGA SoCs provide programmable logic composed of configurable logic blocks, Adaptive Logic Modules (ALMs), lookup tables (LUTs), flip-flops, etc. resources that can be interconnected to construct customized hardware modules or soft-core CPUs.
Cortex-A CCA platforms, we treat the Normal world and all Secure-world software, including the TEE OS, partition manager, and drivers, as untrusted, and trust only the Root-world monitor and the RME hardware.
On Cortex-M MCU systems execute bare-metal firmware or an RTOS with multiple tasks, interrupt handlers, and peripheral-facing components.
We assume the availability of hardware security mechanisms such as secure boot, a TEE, memory protection, and hardware support.
A remote verifier is assumed to maintain the expected hardware/software configuration and verify evidence generated by the embedded device.

Our threat model evolves with the \sys~trust lifecycle.
In the first thrust, during trust establishment, an adversary may attempt to cause an unauthorized hardware or software configuration to be accepted as trusted, for example by leaving a peripheral, interrupt, or DMA master outside the attested domain boundary while still presenting the domain as fully isolated.
In the second thrust, during runtime execution, the adversary may exploit memory-corruption vulnerabilities in applications, libraries, drivers, RTOS components, or interrupt handlers (e.g., ROP, JOP, or data-only attacks) to hijack control flow, corrupt measurement state, or interfere with attestation through interrupts and task transitions.
In the third thrust, a compromised component may attempt to cross compartment boundaries, replay authenticated state, or access unauthorized memory and peripherals, for example through non-control-data/confused-deputy attacks that corrupt shared state without violating control-flow policy or by manipulating CPU-writable DMA descriptor fields to redirect a transfer.
We assume that the root of trust and immutable trusted components are correctly provisioned, secure hardware primitives, and protected key registers cannot be modified by unauthorized software.
Physical invasive attacks and destructive denial-of-service attacks are outside the scope.